\documentclass[12pt,a4paper]{article}

\usepackage[utf8]{inputenc}
\usepackage[T1]{fontenc}
\usepackage[brazil]{babel}
\usepackage{lmodern}
\usepackage{hyperref}

\setlength{\textwidth}{16cm}
\setlength{\textheight}{23cm}
\setlength{\oddsidemargin}{0pt}
\setlength{\evensidemargin}{0pt}
\setlength{\topmargin}{-1cm}
\setlength{\parindent}{0pt}
\setlength{\parskip}{6pt}
\hypersetup{
    colorlinks=true,
    linkcolor=blue,
    urlcolor=blue,
    pdftitle={Relatório Técnico - Chat Multiusuário},
    pdfauthor={Andre Burger},
}

\title{Relatório Técnico\nChat Multiusuário com Sockets, Threads e Tolerância a Falhas}
\author{Andre Burger}
\date{17 de maio de 2026}

\begin{document}
\maketitle

\section{Identificação}

\begin{itemize}
    \item Nome: Andre Burger
    \item Disciplina: Sistemas Distribuídos
    \item Professor: Bruno Dalmazo
    \item Repositório: \url{https://github.com/andre12burger/chat_sd}
\end{itemize}

\section{Motivação e Definição do Problema}

O objetivo do projeto foi desenvolver um chat multiusuário em arquitetura cliente-servidor, acessível via navegador web, com uso explícito de sockets e concorrência baseada em threads. A proposta também exigiu algum mecanismo de tolerância a falhas, de forma que o sistema pudesse ser demonstrado mesmo em caso de indisponibilidade do processo principal.

O problema foi tratado como um caso prático de comunicação interprocesso, concorrência e separação de responsabilidades. O foco esteve no motor do chat, implementado em baixo nível com sockets TCP e threads manuais, enquanto a interface web e a ponte com o navegador foram tratadas em uma camada separada.

\section{Arquitetura do Sistema}

O sistema foi dividido em três partes principais:

\begin{enumerate}
    \item \textbf{Chat Engine} em \texttt{backend/chat\\_engine.py}: servidor TCP puro, responsável por aceitar conexões, criar uma thread por cliente, registrar usuários e fazer broadcast das mensagens.
    \item \textbf{Web Gateway} em \texttt{backend/web\\_gateway.py}: servidor HTTP/WebSocket que entrega a interface web e faz a ponte entre o navegador e o Chat Engine.
    \item \textbf{Frontend} em \texttt{frontend/}: arquivos HTML, CSS e JavaScript separados para a interface do usuário.
\end{enumerate}

A comunicação segue o fluxo: navegador \textrightarrow{} gateway \textrightarrow{} socket TCP \textrightarrow{} engine \textrightarrow{} broadcast. O motor central foi mantido isolado da camada web para facilitar explicação, manutenção e análise do comportamento concorrente.

\section{Regras de Funcionamento}

As regras adotadas foram as seguintes:

\begin{itemize}
    \item múltiplos usuários podem se conectar ao mesmo tempo;
    \item cada conexão de cliente no servidor principal cria uma \texttt{thread} explícita;
    \item o cliente web recebe mensagens em tempo real via Socket.IO;
    \item mensagens enviadas por um usuário passam primeiro pelo servidor central antes de chegar aos demais;
    \item o backend possui mecanismo de tolerância a falhas por meio de um servidor de backup com monitoramento tipo heartbeat.
\end{itemize}

\section{Instalação e Execução}

O projeto utiliza Python 3.10 e as dependências listadas em \texttt{requirements.txt}. Para executar localmente, o fluxo é:

\begin{enumerate}
    \item instalar as dependências com \texttt{pip install -r requirements.txt};
    \item iniciar o \texttt{chat\\_engine.py};
    \item iniciar o \texttt{web\\_gateway.py};
    \item abrir o navegador em \texttt{http://localhost:5001};
    \item testar com múltiplas abas e nomes distintos.
\end{enumerate}

Para facilitar a execução no Windows, foram criados scripts auxiliares como \texttt{run.bat} e \texttt{run.ps1}. Para hospedagem online, a estratégia documentada usa Render.com e UptimeRobot.

\section{Dificuldades Enfrentadas e Soluções Adotadas}

As principais dificuldades foram:

\begin{itemize}
    \item garantir a criação explícita de threads, sem abstrações automáticas que escondessem o requisito;
    \item manter a lista de clientes thread-safe com \texttt{threading.Lock};
    \item separar corretamente frontend, gateway e motor de chat;
    \item preparar o deploy em um ambiente gratuito e estável;
    \item implementar e demonstrar um mecanismo simples de failover.
\end{itemize}

As soluções adotadas foram a divisão em camadas, o uso de locks apenas em seções críticas e a implementação de um servidor de backup por heartbeat. Além disso, a interface foi isolada em arquivos estáticos para manter o código do gateway limpo.

\section{Conclusão e Lições Aprendidas}

O projeto permitiu aplicar conceitos centrais de Sistemas Distribuídos, como comunicação via sockets, sincronização concorrente e organização em múltiplos processos. O uso de threads por cliente no servidor principal deixou clara a diferença entre paralelismo na aplicação e o fluxo de mensagens coordenado pelo servidor.

Como lição principal, ficou evidente que a separação entre responsabilidade de rede, interface e documentação melhora tanto a avaliação acadêmica quanto a manutenção do sistema. A versão final do projeto fornece uma base funcional para demonstração, expansão e escrita do relatório técnico.

\section{Referências do Projeto}

\begin{itemize}
    \item Código-fonte: \url{https://github.com/andre12burger/chat_sd}
    \item Documentação técnica: \texttt{docs/arquitetura.md}
    \item Instruções de execução: \texttt{README.md} e \texttt{docs/SETUP.md}
\end{itemize}

\end{document}
